% Chapter 2, Figure 3: positive components of angular velocity (scan: figures/ch2/fig03.png, PDF 97).
% A true orthographic view (elevation 32 deg, azimuth 45 deg) of the right-handed body axes: TikZ x = E,
% y = F (the rocket's axis, towards the nose), z = D. Positive senses by the right-hand rule: omega_D turns
% E towards F, omega_E turns F towards D, omega_F turns D towards E.
\documentclass[11pt]{standalone}
\usepackage{tamrfig}
\begin{document}
\begin{tikzpicture}[x={(1.15cm,0.61cm)}, y={(-1.15cm,0.61cm)}, z={(0cm,1.38cm)}]
  % rocket: CG at the origin; nose tip at y = 3.7, nose base 2.8, tail 1.65 aft of the CG
  \def\r{0.085}\def\yn{3.7}\def\ynb{2.95}\def\yt{-1.65}
  \def\cr{0.75}\def\ct{0.3}\def\s{0.45}\def\sw{0.45}
  \pgfmathsetmacro\yle{\yt+\cr}\pgfmathsetmacro\ytle{\yle-\sw}\pgfmathsetmacro\ytte{\ytle-\ct}
  \def\phis{48.5} % silhouette angle of the body in this view
  \coordinate (s1t) at ({\r*cos(\phis)},\yt,{\r*sin(\phis)});
  \coordinate (s1n) at ({\r*cos(\phis)},\ynb,{\r*sin(\phis)});
  \coordinate (s2t) at ({-\r*cos(\phis)},\yt,{-\r*sin(\phis)});
  \coordinate (s2n) at ({-\r*cos(\phis)},\ynb,{-\r*sin(\phis)});
  \coordinate (tip) at (0,\yn,0);
  % centre line aft of the tail
  \draw[centerline] (0,\yt,0) -- (0,\yt-1.5,0);
  % fins behind the body: +E and -D
  \draw[outline, fill=white] (\r,\yle,0) -- (\r+\s,\ytle,0) -- (\r+\s,\ytte,0) -- (\r,\yt,0) -- cycle;
  \draw[outline, fill=white] (0,\yle,-\r) -- (0,\ytle,-\r-\s) -- (0,\ytte,-\r-\s) -- (0,\yt,-\r) -- cycle;
  % body and nose
  \def\nose{(s1t) -- (s1n) .. controls ($(s1n)+(0,0.4,0)$) .. (tip) .. controls ($(s2n)+(0,0.4,0)$) .. (s2n) -- (s2t)}
  \fill[white] \nose -- cycle;
  \draw[outline] \nose;
  \draw[outline, canvas is xz plane at y=\ynb] (\phis:\r) arc[start angle=\phis, end angle=\phis+180, radius=\r];
  \draw[outline, fill=white, canvas is xz plane at y=\yt] (0,0) circle[radius=\r];
  % fins in front: +D and -E
  \draw[outline, fill=white] (0,\yle,\r) -- (0,\ytle,\r+\s) -- (0,\ytte,\r+\s) -- (0,\yt,\r) -- cycle;
  \draw[outline, fill=white] (-\r,\yle,0) -- (-\r-\s,\ytle,0) -- (-\r-\s,\ytte,0) -- (-\r,\yt,0) -- cycle;
  % body axes from the centre of gravity
  \draw[thin vec] (0,0,0) -- (0,0,2.3) node[above] {D};
  \draw[thin vec] (0,0,0) -- (2.5,0,0) node[right] {E};
  \draw[thin line] (tip) -- (0,\yn+0.9,0);
  \draw[thin vec] (0,\yn+0.9,0) -- (0,\yn+1.5,0) node[above left=-1pt] {F};
  \node[cg mark] at (0,0,0) {};
  % positive rotations
  \begin{scope}[thin vec]
    \draw[canvas is xy plane at z=1.45] (200:0.55) arc[start angle=200, end angle=500, radius=0.55];
    \draw[canvas is yz plane at x=1.8] (160:0.5) arc[start angle=160, end angle=460, radius=0.5];
    \draw[canvas is xz plane at y=\yn+0.6] (60:0.5) arc[start angle=60, end angle=-240, radius=0.5];
  \end{scope}
  \node at (0.95,0.1,1.45) {$\omega_D$};
  \node at (1.8,-0.55,-0.55) {$\omega_E$};
  \node at (0.55,\yn+0.6,0.5) {$\omega_F$};
\end{tikzpicture}
\end{document}
